% ASSERT_CONVENTION: natural_units=natural, metric_signature=mostly_minus, fourier_convention=physics, coupling_convention=alpha_s, generator_normalization=T(fund)=1/2, covariant_derivative_sign=D_mu=partial_mu-igA
%
% Paper 3: Lattice Signatures of Emergent Supersymmetry at the
%           Conformal Window Boundary
%
% Conventions (Seiberg hep-th/9411149):
%   T(fund) = 1/2 per Weyl fermion
%   b_0 = (11N_c - 2N_f)/3 for N_f Dirac fundamental flavours
%   alpha_s = g_s^2 / (4 pi)
%   R[Q] = (N_f - N_c)/N_f; R[W] = 2
%   N_f counts Dirac flavour pairs (each = 2 Weyl fermions)
%
% EPISTEMIC STATUS: This paper PROPOSES measurements. It does NOT claim
% that SUSY has been observed. All emergent-SUSY predictions carry
% conjectural qualifiers ("if emergent SUSY holds", etc.).
%
\documentclass[11pt,a4paper]{article}

%% ---------- Packages ----------
\usepackage{amsmath,amssymb,amsthm,mathtools}
\usepackage{hyperref}
\usepackage[capitalise,noabbrev]{cleveref}
\usepackage{enumitem}
\usepackage{booktabs}
\usepackage{array}
\usepackage{caption}
\usepackage{xcolor}
\usepackage{slashed}
\usepackage{cite}
\usepackage[margin=2.5cm]{geometry}
\usepackage{multirow}

%% ---------- Hyperref configuration ----------
\hypersetup{
  colorlinks=true,
  linkcolor=blue!70!black,
  citecolor=green!50!black,
  urlcolor=blue!80!black,
}

%% ---------- Theorem environments ----------
\theoremstyle{plain}
\newtheorem{theorem}{Theorem}[section]
\newtheorem{lemma}[theorem]{Lemma}
\newtheorem{proposition}[theorem]{Proposition}
\newtheorem{corollary}[theorem]{Corollary}
\theoremstyle{definition}
\newtheorem{definition}[theorem]{Definition}
\theoremstyle{remark}
\newtheorem{remark}[theorem]{Remark}

%% ---------- Physics macros ----------
\newcommand{\Nc}{N_{\!c}}
\newcommand{\Nf}{N_{\!f}}
\newcommand{\Nfstar}{\Nf^*}
\newcommand{\SU}[1]{\mathrm{SU}(#1)}
\newcommand{\UU}[1]{\mathrm{U}(#1)}
\newcommand{\fund}{\square}
\newcommand{\antifund}{\overline{\square}}
\newcommand{\adj}{\mathrm{adj}}
\newcommand{\Tr}{\mathrm{Tr}}
\newcommand{\ie}{\textit{i.e.}}
\newcommand{\eg}{\textit{e.g.}}
\newcommand{\cf}{\textit{cf.}}
\newcommand{\IR}{\ensuremath{\mathrm{IR}}}
\newcommand{\UV}{\ensuremath{\mathrm{UV}}}
\newcommand{\NSVZ}{\ensuremath{\mathrm{NSVZ}}}
\newcommand{\SQCD}{\ensuremath{\mathrm{SQCD}}}
\newcommand{\QCD}{\ensuremath{\mathrm{QCD}}}
\newcommand{\BZ}{\ensuremath{\mathrm{BZ}}}
\newcommand{\AF}{\ensuremath{\mathrm{AF}}}

%% ---------- Layout ----------
\setlength{\parskip}{0.4em}
\setlength{\parindent}{0em}

%% ================================================================
%%  DOCUMENT
%% ================================================================
\begin{document}

\title{\textbf{Lattice Signatures of Emergent Supersymmetry\\
at the Conformal Window Boundary}}
\author{Notes from the Dualised Standard Model Programme}
\date{March 2026}
\maketitle

\begin{abstract}
Antipin, Mojaza, Pica, and Sannino showed that supersymmetric fixed
points can emerge from non-supersymmetric RG flow in the magnetic dual
of QCD-like theories, at least in the Veneziano limit.  If this
emergent-SUSY phenomenon persists at finite~$\Nc$, it makes specific
predictions for the spectrum and operator dimensions of $\SU{3}$
gauge theories near the lower boundary~$\Nfstar$ of the conformal window.
We identify five lattice-measurable observables that discriminate
between emergent SUSY and ordinary confining or conformal behaviour:
(i)~the scalar-to-pseudoscalar mass ratio $m_\sigma / m_\pi$,
(ii)~the glueball--fermion-bilinear mass relation,
(iii)~the mass anomalous dimension~$\gamma_m$ at the fixed point,
(iv)~a lattice SUSY Ward identity violation parameter, and
(v)~the effective potential for the chiral condensate.  For each
observable, we give the predicted value in the emergent-SUSY scenario
versus the ordinary (non-SUSY) scenario for $\SU{3}$ with
$\Nf = 8$, $10$, and~$12$ fundamental Dirac flavours.
Observables 1 and 3 ($\gamma_m^*$ and $m_\sigma/m_\pi$) are the core diagnostics with direct operator definitions.  Observables 2, 4, and 5 are exploratory extensions contingent on additional operator-mapping work not established here.
Existing lattice data from the LatKMI and LSD collaborations are
compared with these predictions where available.  We propose a
concrete measurement programme to test the emergent-SUSY hypothesis.
All predictions from the emergent-SUSY framework are conjectural.
\end{abstract}

\tableofcontents
\newpage


%% ================================================================
%%  SECTION 1: INTRODUCTION
%% ================================================================
\section{Introduction}
\label{sec:introduction}

The conformal window of a non-abelian gauge theory is the range of
matter content for which the theory possesses an infrared-stable fixed
point: the theory is asymptotically free in the UV but flows to a
conformal field theory in the IR.  For $\SU{\Nc}$ with $\Nf$ Dirac
fundamental flavours, asymptotic freedom is lost at
\begin{equation}\label{eq:AF-bound}
  \Nf^{\AF} = \frac{11\Nc}{2}\,,
\end{equation}
giving $\Nf^{\AF} = 16.5$ for $\SU{3}$.  Below $\Nf^{\AF}$, asymptotic
freedom holds.  As $\Nf$ is decreased, the theory undergoes a transition
at some critical value $\Nfstar$: for $\Nf < \Nfstar$ the theory
confines and breaks chiral symmetry spontaneously, while for
$\Nfstar < \Nf < \Nf^{\AF}$ the theory flows to a conformal fixed point
(the Banks--Zaks fixed point~\cite{BZ1982,Caswell1974}).  The value of
$\Nfstar$ for $\SU{3}$ is not known analytically and is the subject of
active lattice investigation, with current estimates placing it in the
range $8 \lesssim \Nfstar \lesssim 12$~\cite{Hasenbusch2017,
Appelquist2016,Aoki2014,Fodor2012}.

In a remarkable paper, Antipin, Mojaza, Pica, and
Sannino~\cite{AMPS2013} demonstrated that the magnetic dual of a
QCD-like theory, when augmented by Yukawa and quartic couplings,
possesses perturbative fixed points where all couplings simultaneously
take the values predicted by $\mathcal{N}=1$ supersymmetry.  These
SUSY fixed points are reached from generic non-supersymmetric initial
conditions---no fine-tuning is required---and have finite basins of
attraction.  The result was established in the Veneziano limit
($\Nc, \Nf \to \infty$ with $\Nf/\Nc$ fixed), where the fixed-point
couplings are perturbatively small.

If emergent SUSY persists at finite $\Nc$ and governs the physics
near~$\Nfstar$, then the IR dynamics of $\SU{3}$ theories at the
conformal window boundary inherits SUSY predictions.  These predictions
are qualitatively different from those of ordinary confining or
conformal behaviour.  The purpose of this paper is to identify lattice
observables that can distinguish between the two scenarios, derive
the specific SUSY predictions for $\SU{3}$ with $\Nf = 8$, $10$,
and~$12$, and compare with existing lattice data.

\textbf{Epistemic status.}  The emergent-SUSY hypothesis is a
conjecture supported by perturbative evidence in the Veneziano limit.
Whether it survives at $\Nc = 3$ is an open question.  All predictions
labelled ``emergent SUSY'' in this paper are conditional on the
conjecture being correct.  The purpose of our proposed measurements is
precisely to \emph{test} this conjecture, not to assume it.


%% ================================================================
%%  SECTION 2: REVIEW
%% ================================================================
\section{The non-SUSY conformal window and lattice status}
\label{sec:review}

\subsection{Perturbative constraints on the conformal window}

The one-loop coefficient of the beta function for $\SU{\Nc}$ with
$\Nf$ Dirac fundamental flavours is
\begin{equation}\label{eq:b0}
  b_0 = \frac{11\Nc - 2\Nf}{3}\,.
\end{equation}
Asymptotic freedom ($b_0 > 0$) requires $\Nf < 11\Nc/2$.  The
two-loop coefficient is
\begin{equation}\label{eq:b1}
  b_1 = \frac{1}{3}\left(
    34\Nc^2 - \Nf\left(10\Nc + \frac{3(\Nc^2 - 1)}{\Nc}\right)
  \right).
\end{equation}
A perturbative (Banks--Zaks) fixed point exists when $b_0 > 0$ and
$b_1 < 0$, which for $\SU{3}$ gives $\Nf > 8.05$~\cite{BZ1982}.
However, the two-loop estimate for the lower boundary is unreliable
for the actual non-perturbative transition; it merely provides a rough
guide.

\subsection{Non-perturbative estimates}

Several approaches give estimates for $\Nfstar$:

\begin{enumerate}[label=(\roman*)]
  \item \textbf{Schwinger--Dyson (ladder approximation):}
    Chiral symmetry breaking occurs when the mass anomalous dimension
    $\gamma_m$ exceeds unity.  Combined with the two-loop running, this
    gives $\Nfstar \approx 4\Nc = 12$ for $\SU{3}$~\cite{Appelquist1988,
    Cohen1989}.

  \item \textbf{All-orders conjectured beta function
    (Ryttov--Sannino):}  Using a conjectured non-perturbative beta
    function inspired by the NSVZ form, Ryttov and
    Sannino~\cite{RS2008} estimated $\Nfstar \approx (11/4)\Nc = 8.25$
    for $\SU{3}$, with the lower boundary corresponding to
    $\gamma_m = 2$.

  \item \textbf{Functional RG and other methods:}  Various functional
    RG and truncated Schwinger--Dyson studies give
    $\Nfstar \in [8, 12]$~\cite{Gies2006,Braun2010}.
\end{enumerate}

The spread of estimates underscores the need for non-perturbative
lattice calculations.

\subsection{Lattice status for SU(3)}
\label{ssec:lattice-status}

Several collaborations have performed large-scale lattice simulations
of $\SU{3}$ with varying numbers of fundamental flavours:

\textbf{$\boldsymbol{\Nf = 8}$:}  The LatKMI
collaboration~\cite{LatKMI2014_Nf8} and the LSD
collaboration~\cite{LSD2014_Nf8} have studied this theory extensively.
The spectrum shows a light scalar ($\sigma$) state with mass comparable
to the pseudoscalar ($\pi$), reminiscent of a dilaton or
pseudo-Nambu--Goldstone boson of approximate scale invariance.
The mass anomalous dimension is measured to be $\gamma_m \approx
0.4$--$1.0$, depending on the analysis method.  The consensus is that
$\Nf = 8$ is \emph{likely below but near} the conformal window
boundary, with slow (``walking'') running of the coupling.

\textbf{$\boldsymbol{\Nf = 10}$:}  Fewer lattice studies exist at
$\Nf = 10$.  The available data are consistent with either a conformal
fixed point or very slow walking~\cite{Hayakawa2013,Cheng2013}.  This
is the critical region where $\Nfstar$ is expected to lie.

\textbf{$\boldsymbol{\Nf = 12}$:}  This case has generated
significant debate.  The LatKMI collaboration~\cite{LatKMI2016_Nf12}
found evidence for a conformal fixed point with $\gamma_m \approx
0.4$--$0.5$, while Fodor \textit{et al.}~\cite{Fodor2012} argued for
chiral symmetry breaking.  The current consensus leans toward
conformality, with the theory being inside or at the edge of the
conformal window.

\begin{table}[ht]
\centering
\renewcommand{\arraystretch}{1.3}
\caption{Current lattice status for $\SU{3}$ with fundamental Dirac
  flavours.  The mass anomalous dimension $\gamma_m$ is defined via
  $\dim(\bar\psi\psi) = 3 - \gamma_m$.}
\label{tab:lattice-status}
\begin{tabular}{@{} c l c l @{}}
\toprule
$\Nf$ & Status & $\gamma_m$ (lattice) & Key references \\
\midrule
$8$  & Confining/walking & $0.4$--$1.0$ &
  LatKMI~\cite{LatKMI2014_Nf8}, LSD~\cite{LSD2014_Nf8} \\
$10$ & Uncertain & --- &
  \cite{Hayakawa2013,Cheng2013} \\
$12$ & Likely conformal & $0.4$--$0.5$ &
  LatKMI~\cite{LatKMI2016_Nf12}, \cite{Fodor2012} \\
\bottomrule
\end{tabular}
\end{table}


%% ================================================================
%%  SECTION 3: SUSY PREDICTIONS
%% ================================================================
\section{SUSY predictions for the conformal window boundary}
\label{sec:susy-predictions}

In this section we derive the predictions that follow from the
emergent-SUSY hypothesis.  \textbf{All results in this section are
conditional on the conjecture that SUSY governs the IR fixed point
near $\Nfstar$.}

\subsection{The NSVZ beta function and the SUSY conformal window}

In $\mathcal{N}=1$ SQCD with gauge group $\SU{\Nc}$ and $\Nf$
fundamental chiral flavour pairs $(Q_i, \widetilde{Q}_i)$, the exact
NSVZ beta function~\cite{NSVZ1983} is
\begin{equation}\label{eq:NSVZ}
  \beta(\alpha) = -\frac{\alpha^2}{2\pi}\,
  \frac{3\Nc - \Nf(1 - \gamma_Q(\alpha))}
       {1 - \Nc\,\alpha/(2\pi)}\,,
\end{equation}
where $\gamma_Q$ is the anomalous dimension of the fundamental chiral
superfield~$Q$.  The denominator is non-singular in the conformal
window.  The fixed point $\beta(\alpha_*) = 0$ requires
\begin{equation}\label{eq:gammaQ-star}
  \gamma_Q^* = 1 - \frac{3\Nc}{\Nf}\,.
\end{equation}
The SUSY conformal window is the range where the fixed point exists
and respects the unitarity bound $\dim(Q) \geq 1$, \ie\
$\gamma_Q^* \geq -1/2$ (using $\dim(Q) = 1 + \gamma_Q/2$ in the
convention where the free-field dimension is~$1$).  The resulting
window is
\begin{equation}\label{eq:SUSY-window}
  \frac{3}{2}\Nc < \Nf < 3\Nc\,.
\end{equation}
For $\SU{3}$: $4.5 < \Nf < 9$.

\subsection{The mass anomalous dimension}

The key quantity for lattice comparison is the mass anomalous dimension
$\gamma_m$, defined by the scaling dimension of the fermion bilinear:
\begin{equation}\label{eq:gamma-m-def}
  \dim(\bar\psi\psi) = 3 - \gamma_m\,.
\end{equation}
In the SUSY theory, the fermion bilinear is the lowest component of
the meson superfield $M^i{}_j = Q^i \widetilde{Q}_j / \Lambda$.  At a
superconformal fixed point, the dimension of a chiral primary
operator~$\mathcal{O}$ with R-charge $R(\mathcal{O})$ is fixed by the
superconformal algebra~\cite{Seiberg1995,IS1996}:
\begin{equation}\label{eq:dim-R}
  \dim(\mathcal{O}) = \frac{3}{2}\, R(\mathcal{O})\,.
\end{equation}
The meson has R-charge
\begin{equation}\label{eq:R-meson}
  R(M) = 2\,R(Q) = \frac{2(\Nf - \Nc)}{\Nf}\,,
\end{equation}
so its scaling dimension at the fixed point is
\begin{equation}\label{eq:dim-meson}
  \dim(M) = \frac{3(\Nf - \Nc)}{\Nf}\,.
\end{equation}
Identifying $\dim(M) = \dim(\bar\psi\psi) = 3 - \gamma_m$, the SUSY
prediction for the mass anomalous dimension at the fixed point is
\begin{equation}\label{eq:gamma-SUSY}
  \boxed{
    \gamma_m^{\mathrm{SUSY}} = 3 - \frac{3(\Nf - \Nc)}{\Nf}
    = \frac{3\Nc}{\Nf}\,.
  }
\end{equation}

% SELF-CRITIQUE: gamma_m = 3 - 3(N_f - N_c)/N_f = 3 - 3 + 3N_c/N_f
% = 3N_c/N_f.  Sign check OK.  Dimensionless: OK.
% Limits: N_f -> 3N_c => gamma_m -> 1 (unitarity bound from below);
% N_f -> (3/2)N_c => gamma_m -> 2 (unitarity bound from above). OK.

\begin{remark}[Consistency check: conformal window boundaries]
  At the upper boundary $\Nf = 3\Nc$: $\gamma_m^{\mathrm{SUSY}} = 1$.
  The fixed point merges with free-field behaviour.  At the lower
  boundary $\Nf = \frac{3}{2}\Nc$:
  $\gamma_m^{\mathrm{SUSY}} = 2$, which saturates the unitarity bound
  $\dim(\bar\psi\psi) \geq 1$, \ie\ $\gamma_m \leq 2$.  Below this
  point the theory confines.  Both limiting behaviours are correct.
\end{remark}

\subsection{Numerical predictions for SU(3)}
\label{ssec:gamma-numerical}

The emergent-SUSY hypothesis asserts that, near $\Nfstar$ in the
non-SUSY theory, the IR fixed point acquires SUSY-like properties and
the mass anomalous dimension approaches the SUSY value
\eqref{eq:gamma-SUSY}.  For $\SU{3}$:
\begin{equation}\label{eq:gamma-values}
  \gamma_m^{\mathrm{SUSY}}(\Nf) = \frac{9}{\Nf}
  \qquad \Longrightarrow \qquad
  \gamma_m^{\mathrm{SUSY}} =
  \begin{cases}
    9/8  = 1.125 & (\Nf = 8) \\
    9/10 = 0.900 & (\Nf = 10) \\
    9/12 = 0.750 & (\Nf = 12)
  \end{cases}\,.
\end{equation}
These should be compared with the non-SUSY expectation.  In the
ladder Schwinger--Dyson approximation, the lower boundary corresponds
to $\gamma_m = 1$; in the Ryttov--Sannino all-orders beta function, it
corresponds to $\gamma_m = 2$~\cite{RS2008}.  The SUSY values are
distinct from both.  Specifically:

\begin{table}[ht]
\centering
\renewcommand{\arraystretch}{1.3}
\caption{Predicted mass anomalous dimension $\gamma_m$ at the fixed
  point for $\SU{3}$, under different hypotheses for the conformal
  window boundary.  ``BZ'' denotes the perturbative Banks--Zaks
  prediction using the two-loop beta function.  The SUSY values
  use~\eqref{eq:gamma-SUSY}; they apply only \emph{if} emergent SUSY
  governs the fixed point.  The Schwinger--Dyson (SD) value marks the
  onset of chiral symmetry breaking.}
\label{tab:gamma-predictions}
\begin{tabular}{@{} c c c c c @{}}
\toprule
$\Nf$ & $\gamma_m^{\mathrm{SUSY}}$ &
  $\gamma_m^{\BZ}$ (2-loop) &
  $\gamma_m^{\mathrm{SD}}$ &
  $\gamma_m^{\mathrm{lattice}}$ \\
\midrule
$8$  & $1.125$ & $0.46$ & $\sim 1$ &
  $0.4$--$1.0$ \\
$10$ & $0.900$ & $0.26$ & $\sim 1$ & --- \\
$12$ & $0.750$ & $0.13$ & $< 1$ (conformal) &
  $0.4$--$0.5$ \\
\bottomrule
\end{tabular}
\end{table}

The two-loop Banks--Zaks values are much smaller than the SUSY
predictions for all three cases.  The Schwinger--Dyson prediction of
$\gamma_m \to 1$ at $\Nfstar$ differs from the SUSY value of $9/\Nfstar$.
For $\Nfstar = 10$, the distinction is $\gamma_m = 1$ (SD)
versus $\gamma_m = 0.9$ (SUSY)---a~10\% difference that would require
high-precision lattice measurements.  For $\Nfstar = 8$, the SUSY
value $\gamma_m = 1.125$ exceeds the SD value of~$1$, and the
unitarity bound $\gamma_m \leq 2$ is far from saturated.

\subsection{The SUSY spectrum: boson--fermion degeneracies}
\label{ssec:spectrum}

At a superconformal fixed point, operators organise into
superconformal multiplets.  The scalar meson $\sigma$ (the real part
of the chiral condensate) and the pseudoscalar meson $\pi$ (the
imaginary part) are related by a superconformal transformation to the
fermionic component of the same superfield.  The key spectral
predictions are:

\begin{enumerate}[label=(\roman*)]
  \item \textbf{Scalar--pseudoscalar degeneracy:}  If the $\sigma$ and
    $\pi$ mesons belong to the same superconformal multiplet, they must
    have the same scaling dimension.  In a confining theory deformed
    slightly away from the fixed point (to give the states a mass),
    this implies~\cite{Sannino2009_spectrum}
    \begin{equation}\label{eq:sigma-pi}
      \frac{m_\sigma}{m_\pi} \xrightarrow{\text{emergent SUSY}} 1\,.
    \end{equation}
    In contrast, in ordinary QCD-like confining theories,
    $m_\sigma / m_\pi \gg 1$.  In real QCD ($\Nc = 3$, $\Nf = 2$),
    $m_{f_0(500)} / m_\pi \approx 3.6$.

  \item \textbf{Glueball--gluinoball degeneracy:}  In $\mathcal{N}=1$
    super Yang--Mills theory, the glueball (a $0^{++}$ state built from
    $F_{\mu\nu}F^{\mu\nu}$) and the gluinoball (a fermionic state built
    from $\lambda\lambda$, where $\lambda$ is the gluino) belong to the
    same supermultiplet.  If emergent SUSY holds in a theory with an
    adjoint Weyl fermion (as in Sannino's magnetic dual), the glueball
    and the adjoint-fermion bilinear state should become degenerate:
    \begin{equation}\label{eq:glueball}
      \frac{m_{0^{++}}}{m_{\lambda\lambda}}
      \xrightarrow{\text{emergent SUSY}} 1\,.
    \end{equation}
\end{enumerate}


%% ================================================================
%%  SECTION 4: THE FIVE DISCRIMINATING OBSERVABLES
%% ================================================================
\section{Proposed lattice observables}
\label{sec:observables}

We now identify five lattice-measurable quantities that discriminate
between emergent SUSY, ordinary confinement, and conformal behaviour
without emergent SUSY.  For each observable, we give the
prediction in each scenario.

\subsection{Observable 1: scalar-to-pseudoscalar mass ratio
  $m_\sigma / m_\pi$}
\label{ssec:obs1}

The lightest scalar ($\sigma$, quantum numbers $0^{++}$) and
pseudoscalar ($\pi$, quantum numbers $0^{-+}$) mesons are the most
directly accessible hadronic states in lattice simulations.

\textbf{Confining (below window):}  Chiral symmetry is spontaneously
broken.  The pseudoscalar is a (pseudo-)Nambu--Goldstone boson with
$m_\pi \ll m_\sigma$.  In the chiral limit:
\begin{equation}\label{eq:conf-ratio}
  \frac{m_\sigma}{m_\pi} \gg 1
  \qquad\text{(confining)}\,.
\end{equation}
In QCD ($\Nf = 2$): $m_\sigma / m_\pi \approx 3.6$.  For walking
theories ($\Nf$ just below $\Nfstar$), the ratio is expected to
decrease but remain significantly above unity.

\textbf{Conformal without emergent SUSY:}  At a non-SUSY conformal
fixed point, the $\sigma$ and $\pi$ operators have scaling dimensions
determined by the non-SUSY fixed-point dynamics.  In general,
$\dim(\sigma) \neq \dim(\pi)$ because there is no symmetry relating
scalar and pseudoscalar bilinears at a non-SUSY fixed point (they have
different parity and are not related by a SUSY transformation).  Thus:
\begin{equation}\label{eq:BZ-ratio}
  \frac{m_\sigma}{m_\pi} \neq 1
  \qquad\text{(non-SUSY conformal, generically)}\,.
\end{equation}

\textbf{Emergent SUSY:}  At a superconformal fixed point, the scalar
$\sigma$ and pseudoscalar $\pi$ mesons arise as the real and imaginary
parts of the lowest component of the meson chiral superfield $M$.
Within the superconformal multiplet, both have the same scaling
dimension.  Upon deforming away from the fixed point (by a relevant
perturbation that gives the theory a mass gap), the mass degeneracy
is preserved up to corrections proportional to the deformation
parameter:
\begin{equation}\label{eq:SUSY-ratio}
  \boxed{
    \frac{m_\sigma}{m_\pi}
    = 1 + \mathcal{O}\!\left(\frac{m}{m_*}\right)
    \qquad\text{(emergent SUSY)}\,,
  }
\end{equation}
where $m$ is the fermion mass (the deformation) and $m_*$ is the
characteristic scale of the fixed point.

\textbf{Lattice measurement:}  Compute the ratio $m_\sigma / m_\pi$
as a function of the bare fermion mass $m_q$, and extrapolate toward
the chiral limit $m_q \to 0$.  If the ratio approaches unity in the
chiral limit, this is evidence for emergent SUSY.  If it diverges, the
theory confines conventionally.

\textbf{Existing data:}  The LatKMI collaboration~\cite{LatKMI2014_Nf8}
reported $m_\sigma / m_\pi \approx 1.1$--$1.5$ for $\SU{3}$ with
$\Nf = 8$ at the lightest available fermion masses.  This is
suggestively close to the emergent-SUSY prediction of unity, and would disfavour
the ordinary confining prediction $m_\sigma/m_\pi \gg 1$ if confirmed at lighter masses.  However, the
light scalar could also be interpreted as a dilaton (pseudo-Goldstone
boson of spontaneously broken scale invariance) rather than a SUSY
partner of the pion.


\subsection{Observable 2: glueball--fermion-bilinear mass relation}
\label{ssec:obs2}

In a theory containing an adjoint Weyl fermion $\lambda$ in addition
to fundamental fermions, the glueball ($0^{++}$) and the
adjoint-fermion bilinear ($\lambda\lambda$ bound state) should become
degenerate in the emergent-SUSY scenario.

For a non-SUSY $\SU{3}$ theory with $\Nf$ fundamental flavours
only (no adjoint fermion), this observable is not directly available.
However, one can construct a proxy: if emergent SUSY holds, the
glueball mass should satisfy relations inherited from SUSY.  In
$\mathcal{N}=1$ SYM, the gluino condensate sets the scale, and
the glueball mass satisfies~\cite{VY1982_spectrum}
\begin{equation}\label{eq:glueball-condensate}
  m_{0^{++}} \sim \left|\langle\lambda\lambda\rangle\right|^{1/3}\,.
\end{equation}

\textbf{Confining:}  The glueball mass and the fermion-bilinear
condensate are set by the confinement scale $\Lambda$, but there is
no prediction relating $m_{0^{++}}$ to $m_\sigma$.  Generically:
\begin{equation}
  \frac{m_{0^{++}}}{m_\sigma} \sim \mathcal{O}(1)\text{--}\mathcal{O}(10)
  \qquad\text{(confining, no specific prediction)}\,.
\end{equation}

\textbf{Emergent SUSY:}  With an adjoint fermion, the glueball and
gluinoball become degenerate (Eq.~\eqref{eq:glueball}).  Without
the adjoint, the prediction is weaker: the glueball mass should be
related to the fermion-bilinear condensate
via~\eqref{eq:glueball-condensate}, and the glueball should enter
the same tower as the meson states with calculable mass ratios.

\textbf{Lattice measurement:}  Extract the glueball mass from
Wilson-loop or gluonic correlator measurements.  Compare with the
meson spectrum.  A clear separation $m_{0^{++}} \gg m_\sigma$ would
disfavour emergent SUSY; approximate degeneracy would support it.


\subsection{Observable 3: mass anomalous dimension $\gamma_m$}
\label{ssec:obs3}

The mass anomalous dimension $\gamma_m$ is the most precisely
measurable quantity at or near a conformal fixed point.  The
predictions are given in~\cref{tab:gamma-predictions} and repeated
here for emphasis:

\textbf{Emergent SUSY:}
\begin{equation}\label{eq:obs3-SUSY}
  \boxed{
    \gamma_m = \frac{3\Nc}{\Nf}
    = \frac{9}{\Nf}\quad (\SU{3})
  }
\end{equation}

\textbf{Non-SUSY conformal (Banks--Zaks):}  The two-loop estimate gives
values significantly smaller than the SUSY prediction for all $\Nf$
(see~\cref{tab:gamma-predictions}).

\textbf{Schwinger--Dyson:}  The onset of chiral symmetry breaking
occurs at $\gamma_m = 1$.

The discriminating power depends on the location of $\Nfstar$.  If
$\Nfstar \approx 10$, the SUSY value $\gamma_m = 0.9$ is close to
the SD value of~$1$, requiring $\sim 10\%$ precision.  If
$\Nfstar \approx 12$, the SUSY value is $\gamma_m = 0.75$ while the
SD value remains~$1$---a~25\% difference that is accessible to
current lattice methods.

\textbf{Lattice measurement:}  Extract $\gamma_m$ from the scaling of
the Dirac mode number
$\nu(\lambda) \sim \lambda^{3/(1+\gamma_m)}$~\cite{Patella2012}, from
the mass dependence of the pseudoscalar decay constant
$f_\pi \sim m_q^{1/(1+\gamma_m)}$, or from finite-size
scaling~\cite{DeGrand2015}.

\textbf{Existing data:}  For $\Nf = 12$: LatKMI reports $\gamma_m
\approx 0.4$--$0.5$~\cite{LatKMI2016_Nf12}.  The SUSY prediction is
$0.75$.  The discrepancy could indicate that either (a)~emergent SUSY
does not hold at $\Nf = 12$, or (b)~the lattice measurements have
not yet reached the deep IR regime where the fixed-point value is
attained (finite-volume and finite-mass effects can suppress the
effective $\gamma_m$), or (c)~the theory is not at a fixed point at
all.


\subsection{Observable 4: SUSY Ward identity violation}
\label{ssec:obs4}

Even without fundamental SUSY, if a theory flows to a superconformal
fixed point in the IR, the SUSY Ward identities should be approximately
satisfied for low-energy correlation functions.  We propose measuring
the violation of a specific SUSY Ward identity on the lattice.

\subsubsection{Construction of the Ward identity}

In $\mathcal{N}=1$ SUSY, the supercurrent $S_\mu^\alpha$ satisfies
\begin{equation}\label{eq:WI}
  \partial^\mu S_\mu^\alpha = 0
  \qquad\text{(on-shell, in the SUSY limit)}\,.
\end{equation}
When SUSY is broken (or absent), $\partial^\mu S_\mu \neq 0$.  In a
lattice theory, one can construct a discretised version of the
supercurrent and measure its divergence.  This approach has been
successfully applied in lattice studies of $\mathcal{N}=1$ super
Yang--Mills~\cite{DESY2008,Munster2014}.

Define the Ward identity violation parameter:
\begin{equation}\label{eq:WI-violation}
  \boxed{
    \Delta_{\mathrm{SUSY}}(L) \;\equiv\;
    \frac{
      \left\langle\left|\partial^\mu S_\mu(x)\right|^2\right\rangle
    }{
      \left\langle\left|S_\mu(x)\right|^2\right\rangle
    }\,,
  }
\end{equation}
where $L$ is the lattice extent and the expectation values are computed
in the path integral.

\textbf{Confining:}  SUSY is maximally broken.
$\Delta_{\mathrm{SUSY}} = \mathcal{O}(1)$ and does not decrease with
increasing~$L$.

\textbf{Conformal without emergent SUSY:}  The supercurrent is not a
conserved current.  $\Delta_{\mathrm{SUSY}} = \mathcal{O}(1)$ at the
fixed point.

\textbf{Emergent SUSY:}  As the theory flows to the superconformal
fixed point, $\Delta_{\mathrm{SUSY}} \to 0$ in the deep IR.  On a
finite lattice:
\begin{equation}\label{eq:WI-scaling}
  \Delta_{\mathrm{SUSY}}(L) \sim L^{-\Delta}
  \qquad\text{(emergent SUSY)}\,,
\end{equation}
where $\Delta > 0$ is related to the dimension of the leading
SUSY-breaking operator in the fixed-point theory.

\textbf{Lattice measurement:}  The challenge is constructing the
supercurrent $S_\mu$ from the lattice fields.  For a theory with
fundamental fermions and scalars (as in the magnetic dual with
mesons), the supercurrent can be constructed from the fermion and
scalar fields~\cite{DESY2008}.  For a theory with fundamental fermions
only (the electric theory), one must work with composite operators.
We propose using the lattice definition of~\cite{Munster2014},
adapted to include fundamental matter, and measuring
$\Delta_{\mathrm{SUSY}}$ as a function of the lattice extent~$L$ and
the bare fermion mass $m_q$.  A power-law decrease of
$\Delta_{\mathrm{SUSY}}$ with $L$ (or with $1/m_q$ toward the chiral
limit) would be strong evidence for emergent SUSY.


\subsection{Observable 5: effective potential shape and flat directions}
\label{ssec:obs5}

SUSY theories possess moduli spaces of degenerate vacua---flat
directions in the scalar potential that are not lifted by quantum
corrections (within the conformal window).  If emergent SUSY holds, the
effective potential for the chiral condensate should develop approximate
flat directions near $\Nfstar$.

Define the chiral condensate as an order parameter:
\begin{equation}\label{eq:chiral-condensate}
  \Sigma \;\equiv\; \langle\bar\psi\psi\rangle\,.
\end{equation}
In the deep IR of a theory near the fixed point, the effective
potential $V_{\mathrm{eff}}(\Sigma)$ encodes the dynamics of chiral
symmetry breaking.

\textbf{Confining:}  The effective potential has a ``Mexican hat''
shape: $V_{\mathrm{eff}}(\Sigma)$ has a minimum at $\Sigma \neq 0$,
corresponding to spontaneous chiral symmetry breaking.  The curvature
at the minimum sets the sigma mass: $m_\sigma^2 \sim
V''_{\mathrm{eff}}(\Sigma_0)$.

\textbf{Conformal without emergent SUSY:}  The effective potential
has a minimum at $\Sigma = 0$ (no chiral symmetry breaking), with
$V_{\mathrm{eff}}(\Sigma) \sim |\Sigma|^{2/(3-\gamma_m)}$ for small
$\Sigma$.  The curvature depends on $\gamma_m$.

\textbf{Emergent SUSY:}  The effective potential develops approximate
flat directions corresponding to the SUSY moduli space.  Near $\Nfstar$,
the potential should be approximately flat for a range of~$\Sigma$:
\begin{equation}\label{eq:flat-direction}
  V_{\mathrm{eff}}(\Sigma) \approx
  V_0 + \epsilon\,f(\Sigma) + \mathcal{O}(\epsilon^2)\,,
\end{equation}
where $\epsilon$ parameterises the deviation from the superconformal
fixed point and $V_0$ is a constant.  The flatness of the potential
implies a parametrically light scalar excitation---the pseudo-modulus---whose mass
scales as
\begin{equation}\label{eq:pseudo-modulus-mass}
  \boxed{
    m_{\mathrm{pseudo\text{-}mod}}^2 \sim \epsilon\,\Lambda^2
    \qquad\text{(emergent SUSY)}\,,
  }
\end{equation}
where $\Lambda$ is the dynamical scale.  In the exact SUSY limit
$\epsilon \to 0$, the pseudo-modulus becomes massless.

\textbf{Lattice measurement:}  The shape of $V_{\mathrm{eff}}(\Sigma)$
can be reconstructed from the probability distribution of the chiral
condensate on the lattice using constraint
effective potential methods~\cite{CEP}.  A flat region in
$V_{\mathrm{eff}}$ would signal emergent flat directions.
Alternatively, measure the sigma mass as a function of the fermion
mass and check whether $m_\sigma^2$ scales linearly with $m_q$ (as
expected for a pseudo-modulus) rather than remaining of order
$\Lambda^2$ (as in ordinary confinement).


%% ================================================================
%%  SECTION 5: PREDICTIONS FOR SU(3)
%% ================================================================
\section{Expected signals for $\SU{3}$ with $\Nf = 8$, $10$, $12$}
\label{sec:predictions}

We now assemble the predictions of the five observables for the three
most-studied cases.  The predictions depend on whether each $\Nf$
value lies inside the conformal window, outside it (confining), or at
the boundary where emergent SUSY is expected.

\subsection{Scenario assignments}

The emergent-SUSY hypothesis predicts that SUSY-like behaviour appears
precisely at $\Nf = \Nfstar$.  Depending on the (unknown) value of
$\Nfstar$, the three cases are:

\begin{table}[ht]
\centering
\renewcommand{\arraystretch}{1.3}
\caption{Scenario assignments for each $\Nf$ value, depending on the
  location of $\Nfstar$.}
\label{tab:scenarios}
\begin{tabular}{@{} c c c c @{}}
\toprule
$\Nfstar$ & $\Nf = 8$ & $\Nf = 10$ & $\Nf = 12$ \\
\midrule
$\sim 8$  & Boundary (SUSY?) & Conformal & Conformal \\
$\sim 10$ & Confining & Boundary (SUSY?) & Conformal \\
$\sim 12$ & Confining & Confining/walking & Boundary (SUSY?) \\
\bottomrule
\end{tabular}
\end{table}

\subsection{Prediction table}

\Cref{tab:predictions} gives the predicted value of each
observable in each scenario.  The SUSY values apply only
\emph{if} emergent SUSY holds \emph{and} $\Nf$ is at the
boundary.

\begin{table}[ht]
\centering
\renewcommand{\arraystretch}{1.35}
\caption{Predicted values of the five discriminating observables for
  $\SU{3}$ with $\Nf = 8$, $10$, $12$.  ``Conf.''\ = confining,
  ``CW'' = inside conformal window, ``SUSY'' = emergent SUSY at
  boundary.  All SUSY values are \emph{conjectural}.}
\label{tab:predictions}
\begin{tabular}{@{} l c c c @{}}
\toprule
\textbf{Observable} & $\Nf = 8$ & $\Nf = 10$ & $\Nf = 12$ \\
\midrule
\multicolumn{4}{@{}l}{\textit{1.\ Scalar-to-pseudoscalar mass ratio
  $m_\sigma/m_\pi$}} \\
\quad Conf. & $\gg 1$ & $\gg 1$ & --- \\
\quad CW (no SUSY) & --- & $\neq 1$ (generic) & $\neq 1$ (generic) \\
\quad SUSY boundary & $\to 1$ & $\to 1$ & $\to 1$ \\
\midrule
\multicolumn{4}{@{}l}{\textit{2.\ Glueball--meson mass relation}} \\
\quad Conf. & $m_{0^{++}} \sim$ few $\times\, m_\sigma$ & same & --- \\
\quad CW (no SUSY) & --- & unrelated & unrelated \\
\quad SUSY boundary & $m_{0^{++}} \approx m_\sigma$ & same & same \\
\midrule
\multicolumn{4}{@{}l}{\textit{3.\ Mass anomalous dimension
  $\gamma_m$}} \\
\quad Conf. & N/A (no FP) & N/A & --- \\
\quad CW (no SUSY, 2-loop) & --- & $0.26$ & $0.13$ \\
\quad SUSY boundary & $1.125$ & $0.900$ & $0.750$ \\
\midrule
\multicolumn{4}{@{}l}{\textit{4.\ SUSY Ward identity violation
  $\Delta_{\mathrm{SUSY}}$}} \\
\quad Conf. & $\mathcal{O}(1)$ & $\mathcal{O}(1)$ & --- \\
\quad CW (no SUSY) & --- & $\mathcal{O}(1)$ & $\mathcal{O}(1)$ \\
\quad SUSY boundary & $\sim L^{-\Delta}$ & same & same \\
\midrule
\multicolumn{4}{@{}l}{\textit{5.\ Effective potential shape}} \\
\quad Conf. & Mexican hat & same & --- \\
\quad CW (no SUSY) & --- & Single min.\ at 0 & Single min.\ at 0 \\
\quad SUSY boundary & Flat direction & same & same \\
\bottomrule
\end{tabular}
\end{table}


%% ================================================================
%%  SECTION 6: COMPARISON WITH EXISTING DATA
%% ================================================================
\section{Comparison with existing lattice data}
\label{sec:comparison}

We now compare the predictions of~\cref{sec:predictions} with
available lattice results.  We focus on $\Nf = 8$ and $\Nf = 12$,
where substantial data exist.

\subsection{$\Nf = 8$: the light scalar}

The most striking existing result is the light scalar ($\sigma$) state
observed by the LatKMI~\cite{LatKMI2014_Nf8} and
LSD~\cite{LSD2014_Nf8,LSD2016_Nf8} collaborations in $\SU{3}$ with
$\Nf = 8$ fundamental flavours.

\textbf{Observable 1 ($m_\sigma / m_\pi$):}
The LatKMI data show $m_\sigma \approx m_\pi$ at the lightest fermion
masses studied, with $m_\sigma / m_\pi \approx 1.0$--$1.5$ depending
on the fermion mass.  The LSD collaboration also finds a light scalar,
though with somewhat larger ratios.  This is qualitatively consistent
with the emergent-SUSY prediction~\eqref{eq:SUSY-ratio}, but is also
consistent with the dilaton hypothesis (a light scalar arising from
spontaneously broken approximate scale invariance).

Discrimination between the dilaton and SUSY hypotheses requires
additional observables: a dilaton is a scalar singlet, while a SUSY
partner of the pion must carry the same flavour quantum numbers and
have a specific relation to the pseudoscalar decay constant.

\textbf{Observable 3 ($\gamma_m$):}
The mass anomalous dimension for $\Nf = 8$ has been estimated at
$\gamma_m \approx 0.4$--$1.0$.  The SUSY prediction is $\gamma_m =
1.125$.  The upper end of the lattice range ($\gamma_m \sim 1$) is
not far from the SUSY value, but the uncertainties are large.  Higher
precision is needed.  If $\Nf = 8$ is below the conformal window
(as the majority view holds), then $\gamma_m$ is not a fixed-point
value but an effective exponent characterising the walking regime, and
the comparison with the SUSY prediction is less direct.

\subsection{$\Nf = 12$: the conformal case}

\textbf{Observable 3 ($\gamma_m$):}
The LatKMI collaboration reports $\gamma_m \approx 0.4$--$0.5$ for
$\Nf = 12$~\cite{LatKMI2016_Nf12}.  The SUSY prediction is
$\gamma_m = 0.75$.  The discrepancy is~$\sim 50\%$, which is
significant.  If $\Nf = 12$ is well inside the conformal window (far
above $\Nfstar$), then the emergent-SUSY mechanism---which is
predicted to operate specifically at $\Nfstar$---need not apply.  The
measured value could reflect ordinary (non-SUSY) conformal dynamics.

However, if $\Nfstar$ is as high as~$12$, the SUSY prediction
$\gamma_m = 0.75$ should be compared directly.  The discrepancy
with the lattice value of $0.4$--$0.5$ would then constitute evidence
against emergent SUSY at $\Nfstar = 12$.

\subsection{$\Nf = 10$: the missing data}

This is the most critical case.  If $\Nfstar \approx 10$, the
emergent-SUSY predictions are sharpest here: $\gamma_m = 0.9$,
$m_\sigma / m_\pi \to 1$, and $\Delta_{\mathrm{SUSY}} \to 0$.
Unfortunately, very little lattice data exist for $\Nf = 10$.  We
identify this as the highest priority for future lattice simulations.


%% ================================================================
%%  SECTION 7: PROPOSED MEASUREMENT PROGRAMME
%% ================================================================
\section{Proposed measurement programme}
\label{sec:programme}

Based on the analysis above, we propose the following prioritised
measurement programme:

\subsection{Priority 1: $\SU{3}$ with $\Nf = 10$}

This is the most likely location of $\Nfstar$ and the most
under-studied case.  We recommend:
\begin{enumerate}[label=(\alph*)]
  \item Full spectrum calculation: pseudoscalar, scalar, vector,
    and axial-vector mesons.  Extract $m_\sigma / m_\pi$ in the chiral
    limit.
  \item Mass anomalous dimension $\gamma_m$ from the Dirac mode number,
    $f_\pi$ scaling, and finite-size scaling.  Compare with
    $\gamma_m^{\mathrm{SUSY}} = 0.9$.
  \item Glueball spectrum: extract $m_{0^{++}}$ and compare with
    $m_\sigma$.
  \item Multiple fermion masses to enable chiral extrapolation and
    to test whether $m_\sigma^2 \sim m_q$ (pseudo-modulus scaling) or
    $m_\sigma^2 \sim m_q^{2/(1+\gamma_m)}$ (conformal scaling).
\end{enumerate}

\subsection{Priority 2: $\SU{3}$ with $\Nf = 8$ (deeper IR)}

Existing data already hint at a light scalar.  Higher-precision
measurements are needed:
\begin{enumerate}[label=(\alph*)]
  \item Lighter fermion masses to approach the chiral limit and
    determine whether $m_\sigma / m_\pi \to 1$ or diverges.
  \item Improved determination of $\gamma_m$ with reduced systematic
    uncertainties.  The SUSY prediction $\gamma_m = 1.125$ is distinct
    from the SD value of~$1$.
  \item Constraint effective potential measurement to probe the shape
    of $V_{\mathrm{eff}}(\Sigma)$ and search for approximate flat
    directions.
\end{enumerate}

\subsection{Priority 3: SUSY Ward identity measurement (any $\Nf$)}

The SUSY Ward identity violation~\eqref{eq:WI-violation} is
technically challenging but provides the most direct test of emergent
SUSY.  We recommend a pilot study:
\begin{enumerate}[label=(\alph*)]
  \item Construct a lattice discretisation of the supercurrent for
    $\SU{3}$ with $\Nf$ fundamental flavours.  The discretisation
    of~\cite{DESY2008,Munster2014} for pure SYM must be extended to
    include fundamental matter.
  \item Measure $\Delta_{\mathrm{SUSY}}$ at multiple lattice volumes
    for $\Nf = 8$ or $\Nf = 10$.
  \item Check for power-law scaling $\Delta_{\mathrm{SUSY}} \sim
    L^{-\Delta}$.
\end{enumerate}

\subsection{Priority 4: comparison across $\Nf$ values}

A definitive test of the emergent-SUSY hypothesis requires comparing
all five observables across $\Nf = 8$, $10$, and~$12$ (and ideally
$\Nf = 6$ as a control).  The emergent-SUSY predictions apply only
at~$\Nfstar$; away from $\Nfstar$, the ordinary confining or
conformal predictions should hold.  Observing a sharp change in
$m_\sigma / m_\pi$, $\gamma_m$, or $\Delta_{\mathrm{SUSY}}$ at a
specific value of $\Nf$ would be strong evidence for a qualitative
change in the IR dynamics at the conformal window boundary.

\subsection{Summary of discriminating signatures}

\begin{table}[ht]
\centering
\renewcommand{\arraystretch}{1.3}
\caption{Summary of the three scenarios and their distinguishing
  signatures.  A single observable suffices to distinguish confining
  from conformal; discriminating emergent SUSY from ordinary conformal
  requires at least two observables measured at the same~$\Nf$.}
\label{tab:summary-signatures}
\begin{tabular}{@{} l c c c @{}}
\toprule
\textbf{Observable} & \textbf{Confining} &
  \textbf{Conformal} & \textbf{Emergent SUSY} \\
\midrule
$m_\sigma / m_\pi$ & $\gg 1$ & $\neq 1$ (generic) & $\to 1$ \\
$m_{0^{++}} / m_\sigma$ & $\mathcal{O}(1)$--$\mathcal{O}(10)$
  & unrelated & $\to 1$ \\
$\gamma_m$ & N/A & $\ll 3\Nc/\Nf$ & $= 3\Nc/\Nf$ \\
$\Delta_{\mathrm{SUSY}}$ & $\mathcal{O}(1)$ & $\mathcal{O}(1)$
  & $\sim L^{-\Delta}$ \\
$V_{\mathrm{eff}}$ shape & Mexican hat & Min.\ at $\Sigma = 0$ &
  Flat direction \\
\bottomrule
\end{tabular}
\end{table}


%% ================================================================
%%  SECTION 8: CONCLUSIONS
%% ================================================================
\section{Conclusions}
\label{sec:conclusions}

We have identified five lattice-measurable observables that can
discriminate between emergent supersymmetry, ordinary confinement, and
non-SUSY conformal behaviour at the lower boundary of the conformal
window for $\SU{3}$ gauge theories with fundamental fermions.

The emergent-SUSY hypothesis of Antipin, Mojaza, Pica, and
Sannino~\cite{AMPS2013} makes specific, quantitative diagnostic proposals: a mass
anomalous dimension $\gamma_m = 9/\Nf$, a scalar-to-pseudoscalar mass
ratio approaching unity, power-law vanishing of SUSY Ward identity
violations, and the emergence of approximate flat directions in the
effective potential.  These predictions are qualitatively different
from both ordinary confining behaviour and non-SUSY conformal
dynamics.

Existing lattice data provide intriguing but inconclusive evidence.
The light scalar observed at $\Nf = 8$ is consistent with emergent SUSY
but also admits a dilaton interpretation.  The mass anomalous dimension
at $\Nf = 12$ falls below the SUSY prediction, but this may reflect
$\Nf = 12$ being well above~$\Nfstar$.  The critical case $\Nf = 10$
remains largely unexplored on the lattice and should be the highest
priority for future simulations.

The most direct test of emergent SUSY---the measurement of a SUSY Ward
identity violation and its scaling with lattice volume---is technically
challenging but feasible with current lattice technology.  A pilot
study at $\Nf = 8$ or $\Nf = 10$ would provide qualitatively new
information not available from spectral measurements alone.

We emphasise that the emergent-SUSY hypothesis remains a conjecture.
The diagnostic proposals in this paper are the most concrete available consequences
of that conjecture, and the proposed measurement programme is designed
to distinguish between the emergent-SUSY and conventional scenarios.  A positive result would be extraordinary:
it would demonstrate that supersymmetry can emerge dynamically in a
non-supersymmetric theory, with profound implications for our
understanding of strongly coupled gauge dynamics and the naturalness
of the Standard Model.


%% ================================================================
%%  ACKNOWLEDGEMENTS
%% ================================================================
\section*{Acknowledgements}

These notes were developed as part of the Dualised Standard Model
lecture programme.  We thank the LatKMI and LSD collaborations for
making their lattice data publicly available.


%% ================================================================
%%  BIBLIOGRAPHY
%% ================================================================
\begin{thebibliography}{99}

\bibitem{BZ1982}
  T.~Banks and A.~Zaks,
  ``On the phase structure of vector-like gauge theories with massless
  fermions,''
  \textit{Nucl.~Phys.} \textbf{B196} (1982) 189.

\bibitem{Caswell1974}
  W.~E.~Caswell,
  ``Asymptotic behavior of non-abelian gauge theories to two-loop
  order,''
  \textit{Phys.~Rev.~Lett.} \textbf{33} (1974) 244.

\bibitem{Hasenbusch2017}
  A.~Hasenfratz and D.~Schaich,
  ``Nonperturbative beta function of twelve-flavor SU(3) gauge
  theory,''
  \textit{JHEP} \textbf{1802} (2018) 132
  [\texttt{arXiv:1610.10004}].

\bibitem{Appelquist2016}
  T.~Appelquist \textit{et al.}\ (LSD Collaboration),
  ``Strongly interacting dynamics and the search for new physics at
  the LHC,''
  \textit{Phys.~Rev.} \textbf{D93} (2016) 114514
  [\texttt{arXiv:1601.04027}].

\bibitem{Aoki2014}
  Y.~Aoki \textit{et al.}\ (LatKMI Collaboration),
  ``Light composite scalar in twelve-flavor QCD on the lattice,''
  \textit{Phys.~Rev.~Lett.} \textbf{111} (2013) 162001
  [\texttt{arXiv:1305.6006}].

\bibitem{Fodor2012}
  Z.~Fodor \textit{et al.},
  ``Twelve massless flavors and three colors below the conformal
  window,''
  \textit{Phys.~Lett.} \textbf{B703} (2011) 348
  [\texttt{arXiv:1104.3124}].

\bibitem{AMPS2013}
  O.~Antipin, M.~Mojaza, C.~Pica, and F.~Sannino,
  ``Magnetic fixed points and emergent supersymmetry,''
  \textit{JHEP} \textbf{1306} (2013) 037
  [\texttt{arXiv:1105.1510}].

\bibitem{NSVZ1983}
  V.~A.~Novikov, M.~A.~Shifman, A.~I.~Vainshtein, and
  V.~I.~Zakharov,
  ``Exact Gell-Mann--Low function of supersymmetric Yang--Mills
  theories from instanton calculus,''
  \textit{Nucl.~Phys.} \textbf{B229} (1983) 381.

\bibitem{Seiberg1995}
  N.~Seiberg,
  ``Electric--magnetic duality in supersymmetric non-Abelian gauge
  theories,''
  \textit{Nucl.~Phys.} \textbf{B435} (1995) 129
  [\texttt{hep-th/9411149}].

\bibitem{IS1996}
  K.~Intriligator and N.~Seiberg,
  ``Lectures on supersymmetric gauge theories and electric--magnetic
  duality,''
  \textit{Nucl.~Phys.~Proc.~Suppl.} \textbf{45BC} (1996) 1
  [\texttt{hep-th/9509066}].

\bibitem{RS2008}
  T.~A.~Ryttov and F.~Sannino,
  ``Supersymmetry inspired QCD beta function,''
  \textit{Phys.~Rev.} \textbf{D78} (2008) 065001
  [\texttt{arXiv:0711.3745}].

\bibitem{Appelquist1988}
  T.~Appelquist, D.~Karabali, and L.~C.~R.~Wijewardhana,
  ``Chiral hierarchies and the flavor changing neutral current problem
  in technicolor,''
  \textit{Phys.~Rev.~Lett.} \textbf{57} (1986) 957.

\bibitem{Cohen1989}
  A.~G.~Cohen and H.~Georgi,
  ``Walking beyond the rainbow,''
  \textit{Nucl.~Phys.} \textbf{B314} (1989) 7.

\bibitem{Gies2006}
  H.~Gies and J.~Jaeckel,
  ``Chiral phase structure of QCD with many flavors,''
  \textit{Eur.~Phys.~J.} \textbf{C46} (2006) 433
  [\texttt{hep-ph/0507171}].

\bibitem{Braun2010}
  J.~Braun and H.~Gies,
  ``Scaling laws near the conformal window of many-flavor QCD,''
  \textit{JHEP} \textbf{1006} (2010) 060
  [\texttt{arXiv:0912.4168}].

\bibitem{LatKMI2014_Nf8}
  Y.~Aoki \textit{et al.}\ (LatKMI Collaboration),
  ``Light composite scalar in eight-flavor QCD on the lattice,''
  \textit{Phys.~Rev.} \textbf{D89} (2014) 111502
  [\texttt{arXiv:1403.5000}].

\bibitem{LSD2014_Nf8}
  T.~Appelquist \textit{et al.}\ (LSD Collaboration),
  ``Lattice simulations with eight flavors of domain wall fermions
  in SU(3) gauge theory,''
  \textit{Phys.~Rev.} \textbf{D90} (2014) 114502
  [\texttt{arXiv:1405.4752}].

\bibitem{LSD2016_Nf8}
  T.~Appelquist \textit{et al.}\ (LSD Collaboration),
  ``Nonperturbative investigations of SU(3) gauge theory with eight
  dynamical flavors,''
  \textit{Phys.~Rev.} \textbf{D99} (2019) 014509
  [\texttt{arXiv:1807.08411}].

\bibitem{Hayakawa2013}
  M.~Hayakawa \textit{et al.},
  ``Running coupling constant and mass anomalous dimension of
  six-flavor SU(3) gauge theory,''
  \textit{Phys.~Rev.} \textbf{D88} (2013) 094504
  [\texttt{arXiv:1307.6997}].

\bibitem{Cheng2013}
  A.~Cheng, A.~Hasenfratz, G.~Petropoulos, and D.~Schaich,
  ``Scale-dependent mass anomalous dimension from Dirac eigenmodes,''
  \textit{JHEP} \textbf{1307} (2013) 061
  [\texttt{arXiv:1301.1355}].

\bibitem{LatKMI2016_Nf12}
  Y.~Aoki \textit{et al.}\ (LatKMI Collaboration),
  ``Lattice study of conformality in twelve-flavor QCD,''
  \textit{Phys.~Rev.} \textbf{D96} (2017) 014508
  [\texttt{arXiv:1610.07011}].

\bibitem{Patella2012}
  A.~Patella,
  ``A precise determination of the psibar-psi anomalous dimension
  in conformal gauge theories,''
  \textit{Phys.~Rev.} \textbf{D86} (2012) 025006
  [\texttt{arXiv:1204.4432}].

\bibitem{DeGrand2015}
  T.~DeGrand,
  ``Lattice tests of beyond Standard Model dynamics,''
  \textit{Rev.~Mod.~Phys.} \textbf{88} (2016) 015001
  [\texttt{arXiv:1510.05018}].

\bibitem{Sannino2009_spectrum}
  F.~Sannino,
  ``Conformal dynamics for TeV physics and cosmology,''
  \textit{Acta Phys.~Polon.} \textbf{B40} (2009) 3533
  [\texttt{arXiv:0911.0931}].

\bibitem{VY1982_spectrum}
  G.~Veneziano and S.~Yankielowicz,
  ``An effective Lagrangian for the pure $\mathcal{N}=1$
  supersymmetric Yang--Mills theory,''
  \textit{Phys.~Lett.} \textbf{B113} (1982) 231.

\bibitem{DESY2008}
  G.~Bergner \textit{et al.}\ (DESY-M\"unster Collaboration),
  ``The gluino-glue particle and finite size effects in
  supersymmetric Yang--Mills theory,''
  \textit{JHEP} \textbf{1401} (2014) 024
  [\texttt{arXiv:0706.2423}].

\bibitem{Munster2014}
  G.~M\"unster and H.~Stuewe,
  ``The mass of the adjoint pion in $\mathcal{N}=1$
  supersymmetric Yang--Mills theory,''
  \textit{JHEP} \textbf{1405} (2014) 034
  [\texttt{arXiv:1402.6616}].

\bibitem{CEP}
  M.~Fukugita and A.~Ukawa,
  ``Langevin simulation including dynamical quark loops,''
  \textit{Phys.~Rev.~Lett.} \textbf{57} (1986) 503.

\end{thebibliography}

\end{document}
